%قبل از شروع مسئله ترجیحاً چیزی ننویسید. اما اگر خواستید بنویسید، در قالب لیستی (فقرات و شمارش) بنویسید.
%برای اضافه کردن نمره، بعدی از شروع‌مسئله، [x] را اضافه کنید و به جای x  نمره مورد نظر را قرار دهید. برای زیرمسئله‌ها هم به همین صورت است.
%\شروع‌مسئله[10]
\شروع‌مسئله

% صورت اصلی مسئله را در این‌جا بنویسید.
فراز که در گراف اوستا شده است، آوازه‌ش به عنوان یک گراف‌شناس معروف در جهان به گوشه همه رسیده است. به همین دلیل، از او دعوت شده است که در شهری عجیب، یک سیمنار علمی برگزار کند. اما بعضی از دانشمندان این شهر در توانایی‌های فراز تردید‌هایی دارند. در این شهر، 
$n$
میدان وجود دارد و بین برخی از میدان‌ها، جاده‌هایی وجود دارد. می‌توان به میدان‌ها به شکل رئوس یک گراف نگاه کرد و جاد‌ه‌ها، یال‌هایی بین رئوس هستند. در این شهر، بین هر دو میدانی، دقیقا یک مسیر وجود دارد. ضمنا جاده‌ها، طول‌های متفاوتی می‌توانند داشته باشند (طول‌ها اعدادی صحیح و مثبت هستند). حال دانشمندان شهر، نقشه‌ی کامل شهر را به همراه 
$k$
جفت از میدان‌ها به فراز داده‌اند. فراز باید در اردر 
$O(n+k)$
$k$ عدد به آن‌ها بدهد که عدد 
$i$
ام، فاصله‌‌ی دو میدانی است که در جفت 
$i$
ام ذکر شده‌اند. در ضمن فاصله‌ی دو میدان، برابر با 
$XOR$
اعداد مسیر‌ ساده‌ی بین دو میدان است. آیا می‌توانید که الگوریتمی مناسب برای فراز طراحی کنید و او را از این چالش به سلامت عبور دهید؟

\solution{
چون بین هر دو راس دقیقا یک مسیر وجود دارد پس گراف ما درخت است‌. در این صورت کافی است یک راس دلخواه را به عنوان ریشه درخت انتخاب کنیم و از آن راس الگوریتم
 dfs
 را انجام دهیم به طوری که آرایه ای به نام
 cost
 به طول
 n
 در نظر میگیریم. مقدار
 cost
 برای راس ریشه را ۰ در نظر میگیریم و در پیمایش در
 dfs
 مقدار راس فرزند را برابر حاصل
 xor
 مقدار
 cost 
پدر با یال بین آن ها در نظر میگیریم. پس از این عملیات برای هر جفت راس کافی است
 cost
 دو راس را
 xor
 کنیم و به حاصل مورد نظر دست یابیم.
}



%%%%%%%%%%%%%%%%%%%%%%%%%%%%%%%%%%%%%%%%%%%%%%%%%%%%%
% برخی دستورات مفید برای تمرین‌ها. %
%%%%%%%%%%%%%%%%%%%%%%%%%%%%%%%%%%%%%%%%%%%%%%%%%%%%%

% فرمول‌های نمونه:
%$f(x) = \Theta(x)$
%$T(x, y) \leq \log(x) + 2^y + \binom{2n}{n}$
% {\tt code\_function}


% لیست موارد بدون شماره‌گذاری:
%\شروع{فقرات}
%    \فقره مورد اول 
%        \شروع{فقرات}
%            \فقره مورد اول داخلی 
%                \شروع{فقرات}
%                    \فقره اولین مورد خیلی داخلی 
%                    \فقره دومین مورد خیلی داخلی 
%                \پایان{فقرات}
%            \فقره مورد دوم داخلی 
%        \پایان{فقرات}
%    \فقره مورد دوم 
%\پایان{فقرات}

% لیست موارد شماره‌گذاری شده:
%\شروع{شمارش}
%    \فقره یک مورد 
%    \فقره یک مورد دیگر 
%    \فقره باز هم یک مورد دیگر
%\پایان{شمارش}

% برای نوشتن یک فرمول چند خطی که یک شماره بخوره. برای حالت‌های دیگر به سایت زیر مراجعه کنید:
% https://www.overleaf.com/learn/latex/Aligning_equations_with_amsmath
%\begin{align} 
%    \begin{split}
%        (x+y)^3 &= (x+y)^2(x+y) \\
%                &= (x^2+2xy+y^2)(x+y) \\
%                &= (x^3+2x^2y+xy^2) + (x^2y+2xy^2+y^3) \\
%                &= x^3+3x^2y+3xy^2+y^3
%    \end{split}                                 
%\end{align}

% نوشتن ماتریس و آرایه دوبعدی
%\begin{align}
%    A = 
%    \begin{bmatrix}
%        A_{11} & A_{21} \\
%        A_{21} & A_{22}
%    \end{bmatrix}
%\end{align}

% اضافه کردن تصویر با شرح. برای اضافه کردن تصویر بدون شرح، فقط دستور درج‌تصویر را استفاده کنید.
%\شروع{شکل}[h!]
%\تنظیم‌ازوسط
%  \درج‌تصویر[پهنا=0.5\پهنای‌متن]{img.jpg}
%\شرح{شرح تصویر}
%\برچسب{تصویر:۱}
%\پایان{شکل}
%\پایان{زیرمسائل}

%کد زیر را ببینید.
%%\lstinputlisting{source_filename.py}
%\begin{latin}
%\begin{lstlisting}
%# Import the modules
%import sys
%import random
%
%ans = True
%
%while ans:
%question = raw_input("Ask the magic 8 ball a question: (press enter to quit) ")
%
%answers = random.randint(1,8)
%
%if question == "":
%sys.exit()
%
%elif answers == 1:
%print "It is certain"
%
%elif answers == 2:
%print "Outlook good"
%
%elif answers == 3:
%print "You may rely on it"
%
%elif answers == 4:
%print "Ask again later"
%
%elif answers == 5:
%print "Concentrate and ask again"
%
%elif answers == 6:
%print "Reply hazy, try again"
%
%elif answers == 7:
%print "My reply is no"
%
%elif answers == 8:
%print "My sources say no"
%\end{lstlisting}
%\end{latin}
